%%%%%%%%%%%%%%%%
%%% gen 46 %%%
%%%%%%%%%%%%%%%%

\Chapter{46}

\Verse{1}
{
	\hE{וַיִּסַּע יִשְׂרָאֵל וְכל אֲשֶׁר לוֹ וַיָּבֹא בְּאֵרָה שָּׁבַע וַיִּזְבַּח זְבָחִים לֵאלֹהֵי אָבִיו יִצְחָק׃}
}
{
	\eE{
		And Israel and everyone he had traveled and came to
		Beer-sheba, and he offered sacrifices to the {\God} of his
		father, Isaac.
	}
}
{
	Israel and all his people come to Seven Wells.

	\aB{
		And they sacrifice to
		(the sacrifice of		% Isaac
			(their father's God) % YWHW
		)%						% Isaac
		's
		God.%
	}%
	\fC{A number of scholars have a headache after that.}
}

\Verse{2}
{
	\hE{וַיֹּאמֶר אֱלֹהִים לְיִשְׂרָאֵל בְּמַרְאֹת הַלַּיְלָה וַיֹּאמֶר יַעֲקֹב יַעֲקֹב וַיֹּאמֶר הִנֵּנִי׃}
}
{
	\eE{
		And {\God} said to Israel in night visions, and He said,
		“Jacob, Jacob.” And he said, “I’m here.”
	}
}
{
	So God\fA{Who supposely renamed him to Israel.}

	shows up in Israel's mind and goes ``Jacob, Jacob.''

	And he's like ``I'm right here.''

%	\footnote{
%		“Jacob, Jacob.” And he said, “I’m here.” This is the second of three
%		times that this pattern occurs. When Abraham is about to sacrifice
%		Isaac, an angel calls, “Abraham! Abraham!” And he said, “I’m here”
%		(22:11). There it is to save a life. Here it is to assure Jacob that he
%		has nothing to fear in Egypt. Both cases are about reunions of
%		patriarchs with their sons. And this forms the background to the third
%		time the pattern will occur. When Moses comes to the burning bush, God
%		says, “Moses, Moses.” And he said, “I’m here” (Exod 3:4). Like Abraham,
%		Moses is told that Abraham’s children are to be saved. Like Jacob, he is
%		taught not to be afraid. Like both, this is about reunion of the people
%		with the land where the patriarchs are (see comment on Gen. 22:1).
%	}
}

\Verse{3}
{
	\hE{וַיֹּאמֶר אָנֹכִי הָאֵל אֱלֹהֵי אָבִיךָ אַל תִּירָא מֵרְדָה מִצְרַיְמָה כִּי לְגוֹי גָּדוֹל אֲשִׂימְךָ שָׁם׃}
}
{
	\eE{
		And He said, “I am {\God}, your father’s {\God}. Don’t be afraid
		of going down to Egypt, because I’ll make you into a big
		nation there.
	}
}
{
	And he's like ``It's me. God. Your father's God. We got along great,
	me and your dad. Anyways don't be scared. I'll make you into a
	country. I'm thinking of calling it Israel.''
}

\Verse{4}
{
	\hE{אָנֹכִי אֵרֵד עִמְּךָ מִצְרַיְמָה וְאָנֹכִי אַעַלְךָ גַם עָלֹה וְיוֹסֵף יָשִׁית יָדוֹ עַל עֵינֶיךָ׃}
}
{
	\eE{
		I shall go down with you to Egypt, and I shall also bring
		you up. And Joseph will set his hand on your eyes.”
	}
}
{
	I'll go down with you to Egypt.

	I'll bring you back up after.

	Joseph will put his hand on your eyes.
}

\Verse{5}
{
	\hE{וַיָּקם יַעֲקֹב מִבְּאֵר שָׁבַע}
	\hJ{וַיִּשְׂאוּ בְנֵי יִשְׂרָאֵל אֶת יַעֲקֹב אֲבִיהֶם וְאֶת טַפָּם וְאֶת נְשֵׁיהֶם בָּעֲגָלוֹת אֲשֶׁר שָׁלַח פַּרְעֹה לָשֵׂאת אֹתוֹ׃}
}
{
	\eE{And Jacob got up from Beer-sheba.}
	\eJ{And the children of Israel carried Jacob, their father,
		and their infants and their wives in the wagons that Pharaoh
		had sent to carry him.}
}
{
	So Jacob gets up from the dream and from Seven Wells.

	He gets up from all eight of those.

	\aB{
	
	Israel's kids carry Jacob their father.

	They put all the people and stuff in the cars Pharaoh lent them.

	}
}

\Verse{6}
{
	\hP{וַיִּקְחוּ אֶת מִקְנֵיהֶם וְאֶת רְכוּשָׁם אֲשֶׁר רָכְשׁוּ בְּאֶרֶץ כְּנַעַן וַיָּבֹאוּ מִצְרָיְמָה יַעֲקֹב וְכל זַרְעוֹ אִתּוֹ׃}
}
{
	\eP{
		And they took their cattle and their property that they
		had acquired in the land of Canaan, and they came to
		Egypt, Jacob and all his seed with him.
	}
}
{
	Their cows and their other stuff too from back home.

	Then they come to Egypt, Jacob and his seed.
}

\Verse{7}
{
	\hP{בָּנָיו וּבְנֵי בָנָיו אִתּוֹ בְּנֹתָיו וּבְנוֹת בָּנָיו וְכל זַרְעוֹ הֵבִיא אִתּוֹ מִצְרָיְמָה׃}
}
{
	\eP{
		He brought his sons and his grandsons with him, his
		daughters and his granddaughters and all his seed with him
		to Egypt.
	}
}
{
	And Egypt was pregnant with the children of Israel.
}

\Verse{8}
{
	\hP{וְאֵלֶּה שְׁמוֹת בְּנֵי יִשְׂרָאֵל הַבָּאִים מִצְרַיְמָה יַעֲקֹב וּבָנָיו בְּכֹר יַעֲקֹב רְאוּבֵן׃}
}
{
	\eP{
		And these are the names of the children of Israel who came
		to Egypt, Jacob and his sons: Jacob’s firstborn was
		Reuben.
	}
}
{
	Here's the name of everyone who came.
}

\Verse{9}
{
	\hP{וּבְנֵי רְאוּבֵן חֲנוֹךְ וּפַלּוּא וְחֶצְרֹן וְכַרְמִי׃}
}
{
	\eP{
		And Reuben’s sons were Hanoch and Pallu and Hezron and
		Carmi.
	}
}
{
	\Table{|p{0.16\linewidth}|p{0.25\linewidth}|p{0.49\linewidth}|}{
		\hline
		\textbf{Proper name}
			& \textbf{Hebrew}
			& \textbf{Possible Meaning} \\
		\hline
		Hanoch
			& \heb{חֲנוֹךְ} (\emph{Ḥănōḵ})
			& “Dedicated” or “initiated,” from \heb{חנך}. \\
		\hline
		Pallu
			& \heb{פַלּוּא} (\emph{Pallûʾ})
			& Uncertain; sometimes connected with \heb{פלא}, “be distinct, extraordinary.” \\
		\hline
		Hezron
			& \heb{חֶצְרֹן} (\emph{Ḥeṣrōn})
			& Probably related to \heb{חצר}, “enclosure, courtyard, settlement.” \\
		\hline
		Carmi
			& \heb{כַרְמִי} (\emph{Karmî})
			& “Vineyard man” or “my vineyard,” from \heb{כרם}, “vineyard.” \\
		\hline
	}

	\aB{Look Son's sons were Trained and Different and Surrounded By A Wall and Garden.}
}

\Verse{10}
{
	\hP{וּבְנֵי שִׁמְעוֹן יְמוּאֵל וְיָמִין וְאֹהַד וְיָכִין וְצֹחַר וְשָׁאוּל בֶּן הַכְּנַעֲנִית׃}
}
{
	\eP{
		And Simeon’s sons were Jemuel and Jamin and Ohad and
		Jachin and Zohar and Saul, the son of a Canaanite woman.
	}
}
{
	\Table{|p{0.16\linewidth}|p{0.25\linewidth}|p{0.49\linewidth}|}{
		\hline
		\textbf{Proper name}
			& \textbf{Hebrew}
			& \textbf{Possible Meaning} \\
		\hline
		Jemuel
			& \heb{יְמוּאֵל} (\emph{Yəmûʾēl})
			& Uncertain; apparently an \heb{אל}-name, traditionally “day of El.” \\
		\hline
		Jamin
			& \heb{יָמִין} (\emph{Yāmîn})
			& “Right hand” or “south,” from \heb{ימין}. \\
		\hline
		Ohad
			& \heb{אֹהַד} (\emph{ʾŌhad})
			& Uncertain; sometimes associated with “unity” or “joining.” \\
		\hline
		Jachin
			& \heb{יָכִין} (\emph{Yāḵîn})
			& “He establishes” or “may he establish,” from \heb{כון}. \\
		\hline
		Zohar
			& \heb{צֹחַר} (\emph{Ṣōḥar})
			& Probably “brightness” or “whiteness,” from \heb{צחר}. \\
		\hline
		Saul
			& \heb{שָׁאוּל} (\emph{Šāʾûl})
			& “Asked for” or “requested,” from \heb{שאל}. \\
		\hline
	}

	\aB{Hearison's sons were}

	\begin{enumerate}
	\item his 1st son Godday
	\item his 2nd ben Jamin
	\item his 3rd one 1
	\item his 4th son Fou\raisebox{-0.1ex}{n}\raisebox{-0.2ex}{d}\raisebox{-0.1ex}{e}r
	\item his 5th s☉n Light
	\item his 6th was Borrowed
	\end{enumerate}

}

\Verse{11}
{
	\hP{וּבְנֵי לֵוִי גֵּרְשׁוֹן קְהָת וּמְרָרִי׃}
}
{
	\eP{And Levi’s sons were Gershon, Kohath, and Merari.}
}
{
	\Table{|p{0.16\linewidth}|p{0.25\linewidth}|p{0.49\linewidth}|}{
		\hline
		\textbf{Proper name}
			& \textbf{Hebrew}
			& \textbf{Possible Meaning} \\
		\hline
		Gershon
			& \heb{גֵּרְשׁוֹן} (\emph{Gēršōn})
			& Probably related to \heb{גרש}, “drive out”; hence “exile” or “expelled one.” \\
		\hline
		Kohath
			& \heb{קְהָת} (\emph{Qəhāṯ})
			& Uncertain; no secure Hebrew etymology. \\
		\hline
		Merari
			& \heb{מְרָרִי} (\emph{Mərārî})
			& Probably related to \heb{מרר}, “be bitter”: “bitter” or “my bitterness.” \\
		\hline
	}

	\aB{Levi was the father of Exile, Wrinkles, and Bitterness: the children of Levi.}

	\aB{These are the children of the Levites by their names as we have named them.}
}

\Verse{12}
{
	\hP{וּבְנֵי יְהוּדָה עֵר וְאוֹנָן וְשֵׁלָה וָפֶרֶץ וָזָרַח וַיָּמת עֵר וְאוֹנָן בְּאֶרֶץ כְּנַעַן וַיִּהְיוּ בְנֵי פֶרֶץ חֶצְרֹן וְחָמוּל׃}
}
{
	\eP{
		And Judah’s sons were Er and Onan and Shelah and Perez and
		Zerah. And Er and Onan died in the land of Canaan, and
		Perez’s sons were Hezron and Hamul.
	}
}
{
	\Table{|p{0.16\linewidth}|p{0.25\linewidth}|p{0.49\linewidth}|}{
		\hline
		\textbf{Proper name}
			& \textbf{Hebrew}
			& \textbf{Possible Meaning} \\
		\hline
		Er
			& \heb{עֵר} (\emph{ʿĒr})
			& Uncertain; often connected with “awake, watchful.” \\
		\hline
		Onan
			& \heb{אוֹנָן} (\emph{ʾÔnān})
			& Probably from \heb{און}, “strength, vigor.” \\
		\hline
		Shelah
			& \heb{שֵׁלָה} (\emph{Šēlāh})
			& Uncertain; several Hebrew derivations have been proposed. \\
		\hline
		Perez
			& \heb{פֶּרֶץ} (\emph{Pereṣ})
			& “Breach” or “breakthrough,” from \heb{פרץ}, “break through.” \\
		\hline
		Zerah
			& \heb{זֶרַח} (\emph{Zeraḥ})
			& “Rising” or “dawning,” from \heb{זרח}, “rise, shine.” \\
		\hline
		Hezron
			& \heb{חֶצְרוֹן} (\emph{Ḥeṣrôn})
			& Probably from \heb{חצר}, “enclosure, courtyard, settlement”:
			  something like “enclosed place” or “settlement.” \\
		\hline
		Hamul
			& \heb{חָמוּל} (\emph{Ḥāmûl})
			& “Spared” or “pitied,” “mercy,” from \heb{חמל}. \\
		\hline
	}
	
	\aB{
	
	Judah had ᗡＡᗺ (who was bad for some reason (no one knows why))

	And Aonn\fB{That's \paleo{אונן}, not \paleo{א}\hR{\paleo{מ}}\paleo{נון}.}

	Shelah\fB{That's \paleo{שלה}, not \paleo{של}\hR{\paleo{מ}}\paleo{ה}} and Perez and Zerah (ahem)

	Okay so Er was bad and Onan came to an outside end\fC{%
		An old northern English and Scottish idiom meaning to come to a complete,
		disastrous, or ruined finish. Often used to describe an object being smashed
		to pieces, a project failing totally, or a person coming to a bad end.
	}

	Now Judah's seed left him and Tamar got his goat and when he didn't deliver
		Tamar%
		\fC{Lit. TMR; Heb. \heb{תמר}}
		said%
		\fC{Lit. TʾMR. Heb. \heb{תאמר}, the 3rd person feminine singular of the verb ``to say.''}
		``whose stick is this''\fC{Lit. /huzdɪkɪzðɪs/ in the original.}
		and Judah was like ``she's more righteous than I''
		and threw himself on the mercy of the court and then he had twins named
		Buster \& Up and then Buster has kids and names them Mercy and Court.
	}
}

\Verse{13}
{
	\hP{וּבְנֵי יִשָּׂשכָר תּוֹלָע וּפֻוָה וְיוֹב וְשִׁמְרֹן׃}
}
{
	\eP{
		And Issachar’s sons were Tola and Puvah and Job and
		Shimron.
	}
}
{
	\Table{|p{0.16\linewidth}|p{0.25\linewidth}|p{0.49\linewidth}|}{
		\hline
		\textbf{Proper name}
			& \textbf{Hebrew}
			& \textbf{Possible Meaning} \\
		\hline
		Tola
			& \heb{תּוֹלָע} (\emph{Tôlāʿ})
			& “Crimson worm/grub,” the source of scarlet dye. \\
		\hline
		Puvah
			& \heb{פֻוָּה} (\emph{Puvvāh})
			& Uncertain; no secure Hebrew etymology. \\
		\hline
		Job
			& \heb{יוֹב} (\emph{Yôḇ})
			& Uncertain; this form has no secure Hebrew derivation. \\
		\hline
		Shimron
			& \heb{שִׁמְרוֹן} (\emph{Šimrōn})
			& Probably related to \heb{שמר}, “guard, watch, keep.” \\
		\hline
	}

	\aB{Issachar's sons are named Redworm and Blast and Hated and Keeper.}
}

\Verse{14}
{
	\hP{וּבְנֵי זְבֻלוּן סֶרֶד וְאֵלוֹן וְיַחְלְאֵל׃}
}
{
	\eP{And Zebulun’s sons were Sered and Elon and Jahleel.}
}
{
	\Table{|p{0.16\linewidth}|p{0.25\linewidth}|p{0.49\linewidth}|}{
		\hline
		\textbf{Proper name}
			& \textbf{Hebrew}
			& \textbf{Possible Meaning} \\
		\hline
		Sered
			& \heb{סֶרֶד} (\emph{Sereḏ})
			& Uncertain; no secure Hebrew etymology. \\
		\hline
		Elon
			& \heb{אֵלוֹן} (\emph{ʾĒlōn})
			& “Oak” or “terebinth,” from \heb{אלון}. \\
		\hline
		Jahleel
			& \heb{יַחְלְאֵל} (\emph{Yaḥləʾēl})
			& Probably “El waits/hopes,” from \heb{יחל} + \heb{אל}. \\
		\hline
	}
	
	\aB{Zebulun's sons are Dried Up, Tree, and Waiting For God.}
}

\Verse{15}
{
	\hP{אֵלֶּה בְּנֵי לֵאָה אֲשֶׁר יָלְדָה לְיַעֲקֹב בְּפַדַּן אֲרָם וְאֵת דִּינָה בִתּוֹ כּל נֶפֶשׁ בָּנָיו וּבְנוֹתָיו שְׁלֹשִׁים וְשָׁלֹשׁ׃}
}
{
	\eP{
		These were the sons of Leah, to whom she gave birth for
		Jacob in Paddan Aram, and Dinah, his daughter, every
		person of his sons and his daughters, thirty-three.
	}
}
{
	\aB{That's just Leah. She had all those sons, and then she had Dinah.}
}

\Verse{16}
{
	\hP{וּבְנֵי גָד צִפְיוֹן וְחַגִּי שׁוּנִי וְאֶצְבֹּן עֵרִי וַאֲרוֹדִי וְאַרְאֵלִי׃}
}
{
	\eP{
		And Gad’s sons were Ziphion and Haggai and Shuni and
		Ezbon, Eri and Arodi and Areli.
	}
}
{
	\Table{|p{0.16\linewidth}|p{0.25\linewidth}|p{0.49\linewidth}|}{
		\hline
		\textbf{Proper name}
			& \textbf{Hebrew}
			& \textbf{Possible Meaning} \\
		\hline
		Ziphion
			& \heb{צִפְיוֹן} (\emph{Ṣifyōn})
			& Probably related to \heb{צפה}, “watch, look out”: “watcher” or “lookout.” \\
		\hline
		Haggi
			& \heb{חַגִּי} (\emph{Ḥaggî})
			& “Festive” or “my feast,” from \heb{חג}, “festival.” \\
		\hline
		Shuni
			& \heb{שׁוּנִי} (\emph{Šûnî})
			& Uncertain; no secure Hebrew etymology. \\
		\hline
		Ezbon
			& \heb{אֶצְבֹּן} (\emph{ʾEṣbōn})
			& Uncertain; no secure Hebrew etymology. \\
		\hline
		Eri
			& \heb{עֵרִי} (\emph{ʿĒrî})
			& Uncertain; perhaps related to “watchful” or “awake.” \\
		\hline
		Arodi
			& \heb{אֲרוֹדִי} (\emph{ʾĂrôḏî})
			& Uncertain; probably a gentilic or clan name. \\
		\hline
		Areli
			& \heb{אַרְאֵלִי} (\emph{ʾArʾēlî})
			& Probably related to \heb{אראל}, “hero/lion of El”; exact derivation uncertain. \\
		\hline
	}

	\aB{Gad's sons are Watcher, Party, Shuni, Ezbon, Awake, Arodi, and Lion of God.}
}

\Verse{17}
{
	\hP{וּבְנֵי אָשֵׁר יִמְנָה וְיִשְׁוָה וְיִשְׁוִי וּבְרִיעָה וְשֶׂרַח אֲחֹתָם וּבְנֵי בְרִיעָה חֶבֶר וּמַלְכִּיאֵל׃}
}
{
	\eP{
		And Asher’s sons were Imnah and Ishvah and Ishvi and
		Beriah, and Serah was their sister. And Beriah’s sons were
		Hever and Malchiel.
	}
}
{
	\Table{|p{0.16\linewidth}|p{0.25\linewidth}|p{0.49\linewidth}|}{
		\hline
		\textbf{Proper name}
			& \textbf{Hebrew}
			& \textbf{Possible Meaning} \\
		\hline
		Imnah
			& \heb{יִמְנָה} (\emph{Yimnāh})
			& Probably related to \heb{ימין}, “right hand.” \\
		\hline
		Ishvah
			& \heb{יִשְׁוָה} (\emph{Yišwāh})
			& Uncertain; perhaps related to \heb{שוה}, “be equal, level.” \\
		\hline
		Ishvi
			& \heb{יִשְׁוִי} (\emph{Yišwî})
			& Uncertain; perhaps related to \heb{שוה}, “be equal, level.” \\
		\hline
		Beriah
			& \heb{בְרִיעָה} (\emph{Bərîʿāh})
			& Uncertain; later biblical wordplay connects it with “misfortune.” \\
		\hline
		Serah
			& \heb{שֶׂרַח} (\emph{Śeraḥ})
			& Uncertain; sometimes connected with “abundance” or “excess.” \\
		\hline
		Heber
			& \heb{חֶבֶר} (\emph{Ḥeḇer})
			& “Companion” or “association,” from \heb{חבר}, “join.” \\
		\hline
		Malchiel
			& \heb{מַלְכִּיאֵל} (\emph{Malkîʾēl})
			& “My king is El,” from \heb{מלך} + \heb{אל}. \\
		\hline
	}

	\aB{Asher's family is Right Hand, Equal, Equal Again, Misfortune, Plenty, Companion, and My King Is God.}
}

\Verse{18}
{
	\hP{אֵלֶּה בְּנֵי זִלְפָּה אֲשֶׁר נָתַן לָבָן לְלֵאָה בִתּוֹ וַתֵּלֶד אֶת אֵלֶּה לְיַעֲקֹב שֵׁשׁ עֶשְׂרֵה נָפֶשׁ׃}
}
{
	\eP{
		These were the sons of Zil-pah, whom Laban gave to Leah,
		his daughter, and she gave birth to these for Jacob,
		sixteen persons.
	}
}
{
}

\Verse{19}
{
	\hP{בְּנֵי רָחֵל אֵשֶׁת יַעֲקֹב יוֹסֵף וּבִנְיָמִן׃}
}
{
	\eP{
		The sons of Rachel, Jacob’s wife, were Joseph and
		Benjamin.
	}
}
{
}

\Verse{20}
{
	\hP{וַיִּוָּלֵד לְיוֹסֵף בְּאֶרֶץ מִצְרַיִם אֲשֶׁר יָלְדָה לּוֹ אָסְנַת בַּת פּוֹטִי פֶרַע כֹּהֵן אֹן אֶת מְנַשֶּׁה וְאֶת אֶפְרָיִם׃}
}
{
	\eP{
		And Manasseh and Ephraim were born to Joseph in the land
		of Egypt, to whom Asenath, daughter of Poti-phera, priest
		of On, gave birth for him.
	}
}
{
	\Table{|p{0.16\linewidth}|p{0.25\linewidth}|p{0.49\linewidth}|}{
		\hline
		\textbf{Proper name}
			& \textbf{Hebrew}
			& \textbf{Possible Meaning} \\
		\hline
		Asenath
			& \heb{אָסְנַת} (\emph{ʾĀsənaṯ})
			& Egyptian name, often understood as “belonging to Neith”; exact derivation is debated. \\
		\hline
		Poti-phera
			& \heb{פּוֹטִי פֶרַע} (\emph{Pôṭî P̄eraʿ})
			& Egyptian, probably “he whom Ra has given” or “gift of Ra.” \\
		\hline
		On
			& \heb{אֹן} (\emph{ʾŌn})
			& Egyptian \emph{Iunu}, later Heliopolis; probably “the Pillars.” \\
		\hline
		Manasseh
			& \heb{מְנַשֶּׁה} (\emph{Mənaššeh})
			& Gen 41:51 connects it with \heb{נשה}, “cause to forget.” \\
		\hline
		Ephraim
			& \heb{אֶפְרָיִם} (\emph{ʾEfrayim})
			& Gen 41:52 connects it with \heb{פרה}, “be fruitful”: roughly “double fruitfulness.” \\
		\hline
	}

	\aB{Joseph married Belonging To Neith, daughter of Gift Of Ra, priest of The Pillars, and had Forgetting and Double Fruit.}
}

\Verse{21}
{
	\hP{וּבְנֵי בִנְיָמִן בֶּלַע וָבֶכֶר וְאַשְׁבֵּל גֵּרָא וְנַעֲמָן אֵחִי וָרֹאשׁ מֻפִּים וְחֻפִּים וָאָרְדְּ׃}
}
{
	\eP{
		And Benjamin’s sons were Bela and Becher and Ashbel, Gera
		and Naaman, Ehi and Rosh, Muppim and Huppim and Ard.
	}
}
{
	\Table{|p{0.16\linewidth}|p{0.25\linewidth}|p{0.49\linewidth}|}{
		\hline
		\textbf{Proper name}
			& \textbf{Hebrew}
			& \textbf{Possible Meaning} \\
		\hline
		Bela
			& \heb{בֶּלַע} (\emph{Belaʿ})
			& Probably related to \heb{בלע}, “swallow”; exact sense as a name is uncertain. \\
		\hline
		Becher
			& \heb{בֶכֶר} (\emph{Ḇeḵer})
			& Probably related to \heb{בכר}, “firstborn.” \\
		\hline
		Ashbel
			& \heb{אַשְׁבֵּל} (\emph{ʾAšbēl})
			& Uncertain; no secure Hebrew etymology. \\
		\hline
		Gera
			& \heb{גֵּרָא} (\emph{Gērāʾ})
			& Uncertain; no secure Hebrew etymology. \\
		\hline
		Naaman
			& \heb{נַעֲמָן} (\emph{Naʿămān})
			& “Pleasant” or “pleasant one,” from \heb{נעם}. \\
		\hline
		Ehi
			& \heb{אֵחִי} (\emph{ʾĒḥî})
			& “My brother,” from \heb{אח}. \\
		\hline
		Rosh
			& \heb{רֹאשׁ} (\emph{Rōʾš})
			& “Head,” “chief,” or “top.” \\
		\hline
		Muppim
			& \heb{מֻפִּים} (\emph{Muppîm})
			& Uncertain; no secure Hebrew etymology. \\
		\hline
		Huppim
			& \heb{חֻפִּים} (\emph{Ḥuppîm})
			& Uncertain; no secure Hebrew etymology. \\
		\hline
		Ard
			& \heb{אָרְדְּ} (\emph{ʾĀrd})
			& Uncertain; no secure Hebrew etymology. \\
		\hline
	}

	\aB{Benjamin's sons are Swallowed, Firstborn, Ashbel, Gera, Pleasant, My Brother, Head, Muppim, Huppim, and Ard.}
}

\Verse{22}
{
	\hP{אֵלֶּה בְּנֵי רָחֵל אֲשֶׁר יֻלַּד לְיַעֲקֹב כּל נֶפֶשׁ אַרְבָּעָה עָשָׂר׃}
}
{
	\eP{
		These were Rachel’s sons, who were born to Jacob. All the
		persons were fourteen.
	}
}
{
}

\Verse{23}
{
	\hP{וּבְנֵי דָן חֻשִׁים׃}
}
{
	\eP{And Dan’s sons: Hushim.}
}
{
	\Table{|p{0.16\linewidth}|p{0.25\linewidth}|p{0.49\linewidth}|}{
		\hline
		\textbf{Proper name}
			& \textbf{Hebrew}
			& \textbf{Possible Meaning} \\
		\hline
		Hushim
			& \heb{חֻשִׁים} (\emph{Ḥušîm})
			& Uncertain; no secure Hebrew etymology. \\
		\hline
	}

	\aB{Dan had one son. His name was Hushim. We don't know what that means either.}
}

\Verse{24}
{
	\hP{וּבְנֵי נַפְתָּלִי יַחְצְאֵל וְגוּנִי וְיֵצֶר וְשִׁלֵּם׃}
}
{
	\eP{
		And Naphtali’s sons were Jahzeel and Guni and Jezer and
		Shillem.
	}
}
{
	\Table{|p{0.16\linewidth}|p{0.25\linewidth}|p{0.49\linewidth}|}{
		\hline
		\textbf{Proper name}
			& \textbf{Hebrew}
			& \textbf{Possible Meaning} \\
		\hline
		Jahzeel
			& \heb{יַחְצְאֵל} (\emph{Yaḥṣəʾēl})
			& “El apportions/divides,” from \heb{חצה} + \heb{אל}. \\
		\hline
		Guni
			& \heb{גּוּנִי} (\emph{Gûnî})
			& Uncertain; no secure Hebrew etymology. \\
		\hline
		Jezer
			& \heb{יֵצֶר} (\emph{Yēṣer})
			& “Form” or “formation,” from \heb{יצר}, “form, fashion.” \\
		\hline
		Shillem
			& \heb{שִׁלֵּם} (\emph{Šillēm})
			& “Repaid” or “made whole,” from \heb{שלם}. \\
		\hline
	}

	\aB{Naphtali's sons are God Divides, Guni, Formed, and Paid Back.}
}

\Verse{25}
{
	\hP{אֵלֶּה בְּנֵי בִלְהָה אֲשֶׁר נָתַן לָבָן לְרָחֵל בִּתּוֹ וַתֵּלֶד אֶת אֵלֶּה לְיַעֲקֹב כּל נֶפֶשׁ שִׁבְעָה׃}
}
{
	\eP{
		These were the sons of Bilhah, whom Laban gave to Rachel,
		his daughter, and she gave birth to these for Jacob. All
		the persons were seven.
	}
}
{
}

\Verse{26}
{
	\hP{כּל הַנֶּפֶשׁ הַבָּאָה לְיַעֲקֹב מִצְרַיְמָה יֹצְאֵי יְרֵכוֹ מִלְּבַד נְשֵׁי בְנֵי יַעֲקֹב כּל נֶפֶשׁ שִׁשִּׁים וָשֵׁשׁ׃}
}
{
	\eP{
		All the persons of Jacob’s who came to Egypt, who came out
		from his thigh, outside of Jacob’s sons’ wives, all the
		persons were sixty-six.
	}
}
{
}

\Verse{27}
{
	\hP{וּבְנֵי יוֹסֵף אֲשֶׁר יֻלַּד לוֹ בְמִצְרַיִם נֶפֶשׁ שְׁנָיִם כּל הַנֶּפֶשׁ לְבֵית יַעֲקֹב הַבָּאָה מִצְרַיְמָה שִׁבְעִים׃}
}
{
	\eP{
		And Joseph’s sons who were born to him in Egypt were two
		persons. All the persons of Jacob’s house who came to
		Egypt were seventy.
	}
}
{
}

\Verse{28}
{
	\hJ{וְאֶת יְהוּדָה שָׁלַח לְפָנָיו אֶל יוֹסֵף לְהוֹרֹת לְפָנָיו גֹּשְׁנָה וַיָּבֹאוּ אַרְצָה גֹּשֶׁן׃}
}
{
	\eJ{
		And he sent Judah ahead of him to Joseph to direct him in
		advance to Goshen. And they came to the land of Goshen.
	}
}
{
	\Table{|p{0.16\linewidth}|p{0.25\linewidth}|p{0.49\linewidth}|}{
		\hline
		\textbf{Proper name}
			& \textbf{Hebrew}
			& \textbf{Possible Meaning} \\
		\hline
		Goshen
			& \heb{גֹּשֶׁן} (\emph{Gōšen})
			& Probably Egyptian in origin; no secure etymology is known. \\
		\hline
	}

	\aB{Judah goes ahead to show them the way to Goshen. We don't know what Goshen means.}
}

\Verse{29}
{
	\hJ{וַיֶּאְסֹר יוֹסֵף מֶרְכַּבְתּוֹ וַיַּעַל לִקְרַאת יִשְׂרָאֵל אָבִיו גֹּשְׁנָה וַיֵּרָא אֵלָיו וַיִּפֹּל עַל צַוָּארָיו וַיֵּבְךְּ עַל צַוָּארָיו עוֹד׃}
}
{
	\eJ{
		And Joseph hitched his chariot and went up to Israel, his
		father, at Goshen. And he appeared to him and fell on his
		neck and wept on his neck a long time.
	}
}
{
}

\Verse{30}
{
	\hJ{וַיֹּאמֶר יִשְׂרָאֵל אֶל יוֹסֵף אָמוּתָה הַפָּעַם אַחֲרֵי רְאוֹתִי אֶת פָּנֶיךָ כִּי עוֹדְךָ חָי׃}
}
{
	\eJ{
		And Israel said to Joseph, “Let me die now after I’ve seen
		your face, that you’re still alive.”
	}
}
{
}

\Verse{31}
{
	\hJ{וַיֹּאמֶר יוֹסֵף אֶל אֶחָיו וְאֶל בֵּית אָבִיו אֶעֱלֶה וְאַגִּידָה לְפַרְעֹה וְאֹמְרָה אֵלָיו אַחַי וּבֵית אָבִי אֲשֶׁר בְּאֶרֶץ כְּנַעַן בָּאוּ אֵלָי׃}
}
{
	\eJ{
		And Joseph said to his brothers and to his father’s house,
		“Let me go up and tell Pharaoh and say to him, ‘My
		brothers and my father’s house that were in the land of
		Canaan have come to me.
	}
}
{
}

\Verse{32}
{
	\hJ{וְהָאֲנָשִׁים רֹעֵי צֹאן כִּי אַנְשֵׁי מִקְנֶה הָיוּ וְצֹאנָם וּבְקָרָם וְכל אֲשֶׁר לָהֶם הֵבִיאוּ׃}
}
{
	\eJ{
		And the people are shepherds, because they were livestock
		people, and they’ve brought their flock and their oxen and
		all they have.’
	}
}
{
}

\Verse{33}
{
	\hJ{וְהָיָה כִּי יִקְרָא לָכֶם פַּרְעֹה וְאָמַר מַה מַּעֲשֵׂיכֶם׃}
}
{
	\eJ{
		And it will be, when Pharaoh will call you and say, ‘What
		is your occupation?’
	}
}
{
}

\Verse{34}
{
	\hJ{וַאֲמַרְתֶּם אַנְשֵׁי מִקְנֶה הָיוּ עֲבָדֶיךָ מִנְּעוּרֵינוּ וְעַד עַתָּה גַּם אֲנַחְנוּ גַּם אֲבֹתֵינוּ בַּעֲבוּר תֵּשְׁבוּ בְּאֶרֶץ גֹּשֶׁן כִּי תוֹעֲבַת מִצְרַיִם כּל רֹעֵה צֹאן׃}
}
{
	\eJ{
		that you’ll say, ‘Your servants were livestock people from
		our youth until now, both we and our fathers,’ so that
		you’ll live in the land of Goshen, because any shepherd is
		an offensive thing to Egypt.”
	}
}
{
}
